\chapter*{译者说明}
\addcontentsline{toc}{chapter}{译者说明}
\markboth{译者说明}{译者说明}

\section*{底本与范围}
本译本以 American Mathematical Society 2025 年出版的 \emph{Proceedings of Symposia in Pure Mathematics} 第 112.2 卷为底本：\emph{The Langlands Program II. Geometrization of the Langlands Correspondence}，Pierre-Henri Chaudouard、Wee Teck Gan、Tasho Kaletha、Yiannis Sakellaridis 编，共 583 个 PDF 页（印刷页 iii–xi，正文 1–567）。译文覆盖该卷全部六篇文章的正文、前言、各篇摘要与参考文献目录；封面、版权页、丛书信息页与书末索引未译。

本卷是 2022 年 IHES Langlands 纲领暑期学校会议录的第二卷（第 3 部分），主题是 Langlands 对应的几何化与范畴化。

\section*{权利声明}
底本为在版图书，著作权归原作者与 American Mathematical Society 所有。本译本是\textbf{演绎作品}，仅供学习与研究使用，不替代原书；本库不代原作者或出版方授权，权利人通知即撤下。按本库登记规则（R3），图书底本的成品 PDF 不随仓库发布，仓库只保存 \LaTeX{} 源码、工程文件、译注与资料卡。

\section*{体例}
\begin{enumerate}[leftmargin=2em,itemsep=0.2em]
\item \textbf{编号}：定理、命题、引理、推论、定义、注、例、公式、节号的编号一律沿用底本，不重新编号。节号在各篇文章内独立从 1 起，与底本一致。正文中的交叉引用（如「定理 3.2.1」「(2.4.3)」）按底本原样给出。
\item \textbf{页锚}：每页正文左侧（切口侧）的小号灰色数字是底本该页的印刷页码，便于逐页对照原文核校。书末未附索引，读者可按页锚回查。
\item \textbf{术语}：以本工程 \texttt{工程文件/work/glossary.tsv} 为准；同一英文术语全书只用一个译名。人名、地名、机构名与数学家的姓氏一般保留原文，不作音译；「Langlands」「Hecke」「Satake」等已成惯例的译法从众。「stack」译「叠」，「moduli stack」译「模叠」，「coherent sheaf」译「凝聚层」，「constructible sheaf」译「可构造层」。
\item \textbf{公式}：所有公式均按底本页面图像逐符号核对后录入 \LaTeX{}。底本排版中的笔误保留原式，另以脚注说明，不径改。
\item \textbf{文献}：参考文献条目保留英文原貌（含作者、刊名、卷期页码、DOI、MR 编号），仅将「References」标题译为「参考文献」。正文中的方括号引用编号与底本一致。
\item \textbf{译注}：以脚注形式给出，标注「译者注：」。译注只用于补充背景、澄清易混术语、指出底本笔误或说明与通行译法的差异，不替代原文论证。
\item \textbf{语言}：力求直译与可读性的平衡，长句按中文习惯拆分，但不删除原文的逻辑层次；避免三类欧化句式（受事主语句滥用、虚化动词堆叠、「就……来说」套语）。英文术语在首次出现时括注原文。
\end{enumerate}

\section*{生成与校订}
本译本的初稿由机器翻译流水线分块产出，随后经术语一致性检查、数学公式逐页核对、编译审计（0 error、0 未定义引用、0 缺字）与人工抽读。全部译注均经回查底本或外部权威资料（nLab、Stacks Project、原始论文）后写入；无法确证者明确标注为「译者按（上下文推断，无外部依据）」。工程流程、疑点清单与逐条决策记录见 \texttt{工程文件/} 目录。

翻译中的错误与不妥之处在所难免，欢迎通过仓库 issue 指出。
